\section{\texorpdfstring{$L^2$}{L2}中的Fourier变换}
\label{sec:fourier-l2}

周期情形的Parseval关系是无限维的勾股定理。全空间上我们也希望有$\int\lvert f\rvert^2=\int\lvert\widehat f\rvert^2$这样的等式，让Fourier变换像酉矩阵一样保持长度。但这个等式属于$L^2$，而上一节的积分定义要求$f\in L^1$。

在$\mathbb R$上两个空间互不包含，反例如下。$s(x)=\sin(\pi x)/(\pi x)$平方可积却不绝对可积，它的不可积性上一节已经算过。另一方面，函数
\[
u(x)=\lvert x\rvert^{-1/2}\mathbf 1_{\{0<\lvert x\rvert<1\}}
\]
绝对可积，$\int u=4$，却不平方可积，因为$\int_{-1}^1\lvert x\rvert^{-1}\,\mathrm dx$发散。

路线与从有理数造实数相同：在交集$L^1\cap L^2$上用原积分定义变换，证明它保持$L^2$范数，再借助稠密性把它延拓到整个$L^2$。在实轴上，$\sqrt2$不能写成有理数，但它是有理数列的极限；一个$L^2$函数也许无法直接积分，但它是$L^1\cap L^2$中函数的极限。

\paragraph{截取逼近}

延拓之前先确认$L^1\cap L^2$在$L^2$中稠密。对$f\in L^2(\mathbb R^d)$，同时截去远处与过大的取值：
\[
f_n=f\,\mathbf 1_{\{\lvert x\rvert\leq n,\ \lvert f(x)\rvert\leq n\}},\qquad n=1,2,\ldots .
\]
$f_n$的支撑含于半径$n$的球中，取值不超过$n$，所以$f_n$可积。又$\lvert f_n\rvert\leq\lvert f\rvert$，所以$f_n\in L^1\cap L^2$。

对几乎每个$x$，$f(x)$有限，当$n\geq\max\{\lvert x\rvert,\lvert f(x)\rvert\}$时$f_n(x)=f(x)$，所以$\lvert f_n-f\rvert^2\to0$几乎处处。又$\lvert f_n-f\rvert^2\leq\lvert f\rvert^2$可积，由控制收敛定理，
\[
\lVert f_n-f\rVert_2\to0.
\]

用一个例子看截断的量。取$g(x)=1/(1+\lvert x\rvert)$，它平方可积（$\int g^2=2$），但$\int g$发散。这时$g\leq1$，只需截去远处，令$g_n=g\mathbf 1_{\{\lvert x\rvert\leq n\}}$，尾部的能量是
\[
\lVert g-g_n\rVert_2^2=2\int_n^\infty\frac{\mathrm dx}{(1+x)^2}=\frac2{1+n}.
\]
$n=1,9,99$时依次为$1$、$0.2$、$0.02$，虽然慢，但趋于零。

同一截取还给出$L^2$中的平移连续性：$y\to0$时$\lVert\tau_yf-f\rVert_2\to0$。论证与引理~\ref{lem:fourier-fejer}的证明相同。$f_n$有界且可积，所以
\[
\lVert\tau_yf_n-f_n\rVert_2^2\leq2n\,\lVert\tau_yf_n-f_n\rVert_1\to0.
\]
而平移保持$L^2$范数，于是
\[
\lVert\tau_yf-f\rVert_2\leq2\lVert f-f_n\rVert_2+\lVert\tau_yf_n-f_n\rVert_2.
\]

\begin{lemma}[$L^2$中的Gauss光滑化]
\label{lem:fourier-gauss-l2}
设$f\in L^2(\mathbb R^d)$，$t>0$。则$f*G_t$对每个$x$有定义，属于$L^2(\mathbb R^d)$，并且$t\to0$时$\lVert f*G_t-f\rVert_2\to0$。若还有$f\in L^1(\mathbb R^d)$，则$f*G_t\in L^1(\mathbb R^d)$且
\begin{equation}
\widehat{f*G_t}(\xi)=\widehat f(\xi)\,G(t\xi).
\label{eq:fourier-gauss-l2-product}
\end{equation}
\end{lemma}

\begin{proof}
由Cauchy--Schwarz不等式，
\[
\int\lvert f(y)\rvert G_t(x-y)\,\mathrm dy\leq\lVert f\rVert_2\lVert G_t\rVert_2,
\]
所以卷积处处有定义。

$G_t(h)\,\mathrm dh$是概率测度，再用一次Cauchy--Schwarz不等式，得到
\[
\lvert(f*G_t)(x)-f(x)\rvert^2\leq\int G_t(h)\lvert f(x-h)-f(x)\rvert^2\,\mathrm dh.
\]
对$x$积分，交换次序，再换元$h=tu$，于是
\[
\lVert f*G_t-f\rVert_2^2\leq\int G(u)\lVert\tau_{tu}f-f\rVert_2^2\,\mathrm du.
\]
被积函数不超过$4\lVert f\rVert_2^2G(u)$，右边有限，所以$f*G_t-f\in L^2$，从而$f*G_t\in L^2$。由平移连续性与控制收敛定理，右边在$t\to0$时趋于零。

若$f\in L^1$，Tonelli定理给出$\lVert f*G_t\rVert_1\leq\lVert f\rVert_1$。再用Fubini定理与换元$z=x-y$，得到
\[
\widehat{f*G_t}(\xi)=\int f(y)e^{-2\pi\mathrm iy\cdot\xi}\Bigl(\int G_t(z)e^{-2\pi\mathrm iz\cdot\xi}\,\mathrm dz\Bigr)\mathrm dy=\widehat f(\xi)\,\widehat{G_t}(\xi).
\]
换元$z=tu$给出$\widehat{G_t}(\xi)=\widehat G(t\xi)=G(t\xi)$。
\end{proof}

式~\eqref{eq:fourier-gauss-l2-product}是第~\ref{sec:fourier-operations}节卷积规则的一个特例，这里要用到它，所以直接验证。

\subsection{等距延拓}
\label{sec:fourier-l2-extension}

\paragraph{交集上的Plancherel恒等式}

先在一个更小的集合上证明能量关系，那里所有积分都绝对收敛。记
\[
D=\{f\in L^1\cap L^2:\widehat f\in L^1\}.
\]
对$f\in D$，$\widehat f$有界且可积，所以也平方可积；由定理~\ref{thm:fourier-l1-inversion}，$f$几乎处处等于$\widehat f$的逆变换。把一个$f$换成反演公式，再交换积分次序，就得到
\[
\begin{aligned}
\int\lvert f\rvert^2\,\mathrm dx
&=\int\overline{f(x)}\Bigl(\int\widehat f(\xi)e^{2\pi\mathrm ix\cdot\xi}\,\mathrm d\xi\Bigr)\mathrm dx\\
&=\int\widehat f(\xi)\,\overline{\int f(x)e^{-2\pi\mathrm ix\cdot\xi}\,\mathrm dx}\,\mathrm d\xi\\
&=\int\lvert\widehat f\rvert^2\,\mathrm d\xi .
\end{aligned}
\]
交换次序的依据是二重积分绝对收敛，其值为$\lVert f\rVert_1\lVert\widehat f\rVert_1<\infty$。

\begin{theorem}[Plancherel恒等式]
\label{thm:fourier-plancherel-l1l2}
若$f\in L^1(\mathbb R^d)\cap L^2(\mathbb R^d)$，则$\widehat f\in L^2(\mathbb R^d)$，并且
\begin{equation}
\lVert\widehat f\rVert_2=\lVert f\rVert_2 .
\label{eq:fourier-plancherel-l1l2}
\end{equation}
\end{theorem}

\begin{proof}
一般的$f$未必属于$D$，用Gauss光滑化把它变成$D$中的元素。令$f_t=f*G_t$。由引理~\ref{lem:fourier-gauss-l2}，$f_t\in L^1\cap L^2$，且$\widehat{f_t}=\widehat f\,G(t\,\cdot)$。$\widehat f$有界、$G(t\,\cdot)$可积，所以$\widehat{f_t}\in L^1$，即$f_t\in D$。

对$f_t$用上面的恒等式，并注意$G(t\xi)^2=e^{-2\pi t^2\lvert\xi\rvert^2}$，得到
\[
\int\lvert\widehat f(\xi)\rvert^2e^{-2\pi t^2\lvert\xi\rvert^2}\,\mathrm d\xi=\lVert f_t\rVert_2^2.
\]
令$t\downarrow0$。左边的被积函数单调增加到$\lvert\widehat f\rvert^2$，由单调收敛定理，左边趋于$\int\lvert\widehat f\rvert^2$。右边由引理~\ref{lem:fourier-gauss-l2}趋于$\lVert f\rVert_2^2$。
\end{proof}

用方框核对。$\lVert B\rVert_2^2=1$，定理说$\int(\sin\pi\xi/(\pi\xi))^2\,\mathrm d\xi=1$。这个积分用初等方法并不好算，下面还会用反演另外验证。

\paragraph{延拓到$L^2$}

现在$f\mapsto\widehat f$是稠密子空间$L^1\cap L^2$到完备空间$L^2$的线性等距映射。第~\ref{chap:metric-space-maps}章证明了：稠密子集上的一致连续映射，若目标空间完备，就有唯一的连续延拓，并且Lipschitz常数保持不变。把它用于这一等距映射，就得到通常所说的\term{Plancherel定理}（Plancherel theorem）。

\begin{theorem}[Plancherel定理]
\label{thm:fourier-plancherel}
存在唯一的连续映射$\mathcal F:L^2(\mathbb R^d)\to L^2(\mathbb R^d)$，使$\mathcal Ff=\widehat f$对$f\in L^1\cap L^2$成立。$\mathcal F$是线性等距：$\lVert\mathcal Ff\rVert_2=\lVert f\rVert_2$。对任意$f\in L^2$，下式在$L^2$范数下成立：
\begin{equation}
\mathcal Ff=\lim_{n\to\infty}\int_{\lvert x\rvert\leq n}f(x)e^{-2\pi\mathrm ix\cdot(\cdot)}\,\mathrm dx .
\label{eq:fourier-plancherel-limit}
\end{equation}
\end{theorem}

\begin{proof}
存在性与唯一性由上述延拓定理给出。

等距性：对$f\in L^2$取$g_n\in L^1\cap L^2$，$g_n\to f$，由$\mathcal F$的连续性与式~\eqref{eq:fourier-plancherel-l1l2}，得到
\[
\lVert\mathcal Ff\rVert_2=\lim_n\lVert\widehat{g_n}\rVert_2=\lim_n\lVert g_n\rVert_2=\lVert f\rVert_2.
\]
线性同样由$L^1\cap L^2$上的线性取极限得到。

最后，$f\mathbf 1_{\{\lvert x\rvert\leq n\}}$在有限测度的球上平方可积，由Cauchy--Schwarz不等式也可积，并在$L^2$中收敛到$f$，所以它们的变换收敛到$\mathcal Ff$。
\end{proof}

此后对$f\in L^2$也写$\widehat f=\mathcal Ff$，称为$f$的Fourier变换；在$L^1\cap L^2$上，这一记号与原积分定义一致。

\paragraph{频域元素是等价类}

延拓得到的$\mathcal Ff$是$L^2$中的元素，即几乎处处相等的函数组成的等价类，不是一个在每个频率处都有确定值的函数。对一般的$f\in L^2$，积分$\int f(x)e^{-2\pi\mathrm ix\cdot\xi}\,\mathrm dx$不必绝对收敛，式~\eqref{eq:fourier-plancherel-limit}的极限也只在范数意义下成立，不保证对某个固定的$\xi$收敛。

例如$s(x)=\sin(\pi x)/(\pi x)$平方可积而不绝对可积，对它不能用定义~\ref{def:fourier-transform-l1}在任何一个频率处读取变换。下一小节将算出$\mathcal Fs=\mathbf 1_{[-1/2,1/2]}$，这是$L^2$中的等式，频率$\pm1/2$处的值没有意义，正如改变方框在端点的取值不改变它所代表的元素。这与$L^1$情形形成对照：那里$\widehat f$是处处有定义的连续函数，可以在单个频率处求值。

\subsection{\texorpdfstring{$L^2$}{L2}反演与Parseval关系}
\label{sec:fourier-l2-inversion}

\paragraph{反演与内积保持}

在$L^2$中，反演有一个干净的形式。记反射$(Rf)(x)=f(-x)$，它是$L^2$上的等距。对$f\in L^1$，直接换元得到
\[
\widehat{Rf}(\xi)=\int f(-x)e^{-2\pi\mathrm ix\cdot\xi}\,\mathrm dx=\widehat f(-\xi).
\]
两边都连续地依赖$f$，所以$\mathcal FR=R\mathcal F$在$L^2$上成立。逆变换式~\eqref{eq:fourier-inverse}恰为$\check h=R\widehat h$，因此它延拓到$L^2$后的映射就是$R\mathcal F$。

\begin{theorem}[$L^2$反演与Parseval关系]
\label{thm:fourier-l2-inversion}
$\mathcal F$是$L^2(\mathbb R^d)$到自身的双射，逆映射为$\mathcal F^{-1}=R\mathcal F$，并且$\mathcal F\mathcal F=R$。对任意$f,g\in L^2(\mathbb R^d)$，
\begin{equation}
\int_{\mathbb R^d}\overline{f(x)}\,g(x)\,\mathrm dx=\int_{\mathbb R^d}\overline{\widehat f(\xi)}\,\widehat g(\xi)\,\mathrm d\xi .
\label{eq:fourier-l2-parseval}
\end{equation}
\end{theorem}

\begin{proof}
先证$R\mathcal F\mathcal F=I$。对$f\in D$，$\widehat f\in L^1\cap L^2$，所以$R\mathcal F(\widehat f)$可以用积分计算，它就是$\widehat f$的逆变换，由定理~\ref{thm:fourier-l1-inversion}等于$f$。所以$R\mathcal F\mathcal F=I$在$D$上成立。

定理~\ref{thm:fourier-plancherel-l1l2}的证明表明，每个$f\in L^1\cap L^2$都是$D$中元素$f*G_t$的$L^2$极限，而$L^1\cap L^2$在$L^2$中稠密，所以$D$在$L^2$中稠密。两边都是连续映射，故$R\mathcal F\mathcal F=I$在$L^2$上成立。

由$R^2=I$得$\mathcal F\mathcal F=R$。再由$\mathcal FR=R\mathcal F$得$\mathcal F(R\mathcal F)=R\mathcal F\mathcal F=I$。所以$\mathcal F$是双射，逆为$R\mathcal F$。

内积保持来自范数保持。对第二个变量线性的复内积，由范数表示的公式是
\[
\langle f,g\rangle=\frac14\sum_{k=0}^3\mathrm i^{-k}\lVert f+\mathrm i^kg\rVert_2^2.
\]
$\mathcal F$线性且保持每个$f+\mathrm i^kg$的范数，所以保持内积。
\end{proof}

这就是有限维酉矩阵的全部性质：保持长度和内积，逆等于“共轭转置”。把$\mathcal F$作用两次等于反射，这一点第~\ref{sec:fourier-discrete}节会用$4\times4$矩阵具体看到。式~\eqref{eq:fourier-l2-parseval}是全空间上的Parseval关系，取$g=f$即能量关系$\lVert\widehat f\rVert_2=\lVert f\rVert_2$，在本章的归一化下两边没有额外常数。

现在可以算出$\mathcal Fs$。方框$B$属于$L^1\cap L^2$，$\widehat B=s$。由$\mathcal F\mathcal F=R$，并注意$B$是偶函数，得到
\[
\mathcal Fs=\mathcal F\mathcal FB=RB=B.
\]
能量关系同时给出
\[
\int_{\mathbb R}\Bigl(\frac{\sin\pi x}{\pi x}\Bigr)^2\,\mathrm dx=\lVert B\rVert_2^2=1.
\]

\paragraph{同一公式的不同含义}

反演公式~\eqref{eq:fourier-l1-inversion}在本章出现了三次，含义各不相同，表~\ref{tab:fourier-inversion-meanings}把它们并列。把积分区域限制在$\lvert\xi\rvert\leq R$得到的截断积分，是否对几乎每个$x$逐点收敛，是深刻得多的问题，本书不讨论。

\begin{BookInlineTable}
  \centering
  \small
  \begin{tabular}{@{}lll@{}}
    \toprule
    条件 & 右边的含义 & 等式的意义\\
    \midrule
    $f,\widehat f\in L^1$ & 绝对收敛的积分 & 几乎处处；连续点处逐点\\
    一维，局部分段$C^1$ & 对称截断的极限 & 每一点；跳跃点取平均\\
    $f\in L^2$ & 截断积分的$L^2$极限 & 范数意义\\
    \bottomrule
  \end{tabular}
  \caption{反演公式在三种条件下的含义。}
  \label{tab:fourier-inversion-meanings}
\end{BookInlineTable}

\paragraph{归一化与常数}

不同文献的Fourier变换只差自变量伸缩与常数因子，但反演与能量公式中的常数必须与定义配套。表~\ref{tab:fourier-normalizations}比较三种常见约定，其中第二、三行用角频率$\omega=2\pi\xi$。

\begin{BookInlineTable}
  \centering
  \small
  \begin{tabularx}{\linewidth}{@{}>{\raggedright\arraybackslash}p{20mm}>{\raggedright\arraybackslash}X>{\raggedright\arraybackslash}p{28mm}@{}}
    \toprule
    约定 & 正变换与反变换 & 能量关系\\
    \midrule
    频率$\xi$（本章） & $\widehat f(\xi)=\int fe^{-2\pi\mathrm ix\cdot\xi}\,\mathrm dx$，$f=\int\widehat fe^{2\pi\mathrm ix\cdot\xi}\,\mathrm d\xi$ & $\int\lvert f\rvert^2=\int\lvert\widehat f\rvert^2$\\
    角频率，正变换不加因子 & $F(\omega)=\int fe^{-\mathrm ix\cdot\omega}\,\mathrm dx$，$f=(2\pi)^{-d}\int Fe^{\mathrm ix\cdot\omega}\,\mathrm d\omega$ & $\int\lvert f\rvert^2=(2\pi)^{-d}\int\lvert F\rvert^2$\\
    角频率，对称归一 & 正、反变换各乘$(2\pi)^{-d/2}$ & $\int\lvert f\rvert^2=\int\lvert F\rvert^2$\\
    \bottomrule
  \end{tabularx}
  \caption{$\mathbb R^d$上Fourier变换的三种常见归一化。各行之间的换算是$F(\omega)=\widehat f(\omega/2\pi)$，对称归一化再乘以$(2\pi)^{-d/2}$；所有结论只差这些常数。}
  \label{tab:fourier-normalizations}
\end{BookInlineTable}

以第二行为例核对常数。作换元$\omega=2\pi\xi$，$\mathrm d\omega=(2\pi)^d\,\mathrm d\xi$，得到
\[
\int\lvert F(\omega)\rvert^2\,\mathrm d\omega=\int\lvert\widehat f(\omega/2\pi)\rvert^2\,\mathrm d\omega=(2\pi)^d\int\lvert\widehat f(\xi)\rvert^2\,\mathrm d\xi,
\]
所以能量关系中出现$(2\pi)^{-d}$。周期情形同样如此：第~\ref{sec:fourier-series}节把$1/(2\pi)$放进内积，Parseval关系才不含常数。
